\chapter{Сравнение индексов Каллиаса и Тёплица по функции склейки}

Предыдущая проверка получила одинаковые числа $+15/-15$ двумя независимыми
способами, но запретил считать численное совпадение доказательством общего
класса. Теперь сравниваются не индексы, а сами граничные символы.

\section{Литературная форма задачи}

Бунке показал, что индекс тёплицева оператора равен индексу связанного
оператора типа Каллиаса; см. \path{arXiv:math/9911177}. Современная
формулировка Дирака--Шрёдингера сводит класс Каллиаса к спариванию
положительного собственного проектора массы на компактной граничной
гиперповерхности с её классом Дирака; см. \path{arXiv:2407.02993}.
Абстрактная версия того же спаривания для спектральных троек дана в
\path{arXiv:2108.06368}.

Прямое применение тёплицевой компрессии Бунке потребовало бы отдельно
проверить спектральную щель и изолированность нуля для выбранного
неограниченного оператора. Эти аналитические гипотезы в проекте пока не
установлены. Поэтому настоящая проверка использует более слабый, но точный
маршрут:
\begin{equation}
 \text{индекс Каллиаса на }\mathbb R^3
 \longrightarrow K^0(S^2_\infty)
 \xrightarrow{\mathrm{clutch}}
 K^1(S^1)
 \xrightarrow{\partial_{\mathcal T}}
 K^0(\mathcal K).
 \label{eq:v6-callias-toeplitz-route}
\end{equation}

\section{Хопфова линия и её экваториальный символ}

Для $n=(\sin\theta\cos\varphi,\sin\theta\sin\varphi,\cos\theta)$
положительный проектор массы спинорного накрытия равен
\begin{equation}
 P_+(n)=\frac{I_2+n\cdot\sigma}{2}.
\end{equation}
На северной и южной картах выберем нормированные собственные векторы
\begin{align}
 v_N(\theta,\varphi)
 &=\begin{pmatrix}\cos(\theta/2)\\
 e^{i\varphi}\sin(\theta/2)\end{pmatrix},\\
 v_S(\theta,\varphi)
 &=\begin{pmatrix}e^{-i\varphi}\cos(\theta/2)\\
 \sin(\theta/2)\end{pmatrix}.
\end{align}
Оба задают $P_+=vv^*$ в своей карте. На экваторе выполнено точное
соотношение
\begin{equation}
 v_N=e^{i\varphi}v_S.
 \label{eq:v6-hopf-clutching-z}
\end{equation}
Следовательно, функция склейки положительной хопфовой линии есть
\begin{equation}
 g_L(z)=z,
 \qquad z=e^{i\varphi},
 \qquad c_1(L)=\operatorname{wind}(g_L)=1.
\end{equation}

Это уже не численная аппроксимация интеграла Черна, а явный представитель
того же класса.

\section{Коэффициентный ранг пятнадцать}

Возьмём тот же канонический коэффициентный проектор Тома V
\begin{equation}
 q_0=p_{15}\otimes P_0\in M_{105}(\mathbb C),
 \qquad \rank q_0=15.
\end{equation}
Функция склейки расслоения $L\otimes q_0$ равна
\begin{equation}
 g_{15}(z)=zq_0+1-q_0.
 \label{eq:v6-callias-clutching-v15}
\end{equation}
Но правая часть \eqref{eq:v6-callias-clutching-v15} дословно совпадает с
первой компонентой явного Real-унитария Тома V:
\begin{equation}
 V_+(z)=zq_0+1-q_0.
\end{equation}
Его определитель на коэффициентном образе равен $z^{15}$, поэтому
\begin{equation}
 \operatorname{wind}\det g_{15}=15.
\end{equation}

Сопряжённая ветвь спинорного накрытия имеет линию $\overline L$, переход
$z^{-1}$ и символ
\begin{equation}
 g_{-15}(z)=z^{-1}\overline q_0+1-\overline q_0=V_-(z).
\end{equation}
Тем самым совпадает и Real-обмен двух ориентаций, а не только абсолютное
целое число.

\begin{theorem}[Тождество граничного класса]
Положительное собственное расслоение массы спинорного накрытия
$n\cdot\sigma$, тензорно умноженное на $q_0$, имеет унитарий склейки,
тождественно равный компоненте $V_+$ Real-символа Тёплица из Тома V.
Сопряжённая ветвь тождественно равна $V_-$. Поэтому граница Каллиаса и
граница Тёплица представляют один и тот же сопряжённый граничный
$K$-класс с коэффициентом пятнадцать.
\end{theorem}

\section{Согласование знаков индекса}

Граничная теорема Каллиаса даёт, с точностью до выбранной ориентации,
\begin{equation}
 \operatorname{ind}\mathcal D_{\rm Callias}
 =\left\langle [L\otimes q_0],[D_{S^2}]\right\rangle
 =15.
\end{equation}
Для стандартной ориентации пространства Харди теорема Тёплица даёт
\begin{equation}
 \operatorname{ind}T_{g_{15}}
 =-\operatorname{wind}\det g_{15}=-15.
\end{equation}
Разность знаков является обычной разностью ориентационных конвенций
границ Каллиаса и сжатия Харди. При одновременном обращении
ориентации получается сопряжённая пара
\begin{equation}
 (+15,-15)_{\rm Callias}
 \longleftrightarrow
 (-15,+15)_{\rm Toeplitz}.
\end{equation}
Абсолютный Real-класс и нормированный вес $15/105=1/7$ совпадают.

\section{Что закрыто и что осталось}

Сравнительный $K$-теоретический мост теперь закрыт: хопфова сфера
пространственного дефекта и окружность Тёплица связаны обычной функцией
склейки на границе двух карт той же сферы. Окружность здесь не является новым
компактным измерением и не отождествляется с внешним множителем $S^1$ в
$\mathbb{RP}^3\times S^1$.

Однако операторное и физическое замыкание ещё не достигнуто. Не построено
унитарное отождествление полного трёхмерного оператора Каллиаса с
оператором Харди; оно и не требуется для равенства индексов. Существеннее,
что конечный родитель всё ещё не вывел комплексный носитель ранга два, на
котором масса $n\cdot\sigma$ действует независимо от пространственного
клиффордова модуля.

Если такой носитель существует, ненулевой индекс Каллиаса уже гарантирует
избыток пятнадцати квадратично интегрируемых нулевых мод одной ориентации.
Но без вывода носителя из $H_{15}/M_{35}$ это остаётся условной
локализационной теоремой, а не рождённой физической материей.

\begin{nogocriterion}[Класс не создаёт отсутствующий носитель]
Тождество символов склейки доказывает равенство граничных $K$-классов,
но не разрешает добавлять пространство спинорного накрытия ранга два вручную.
Физический вывод фермионов закрывается только после построения этого
носителя и его массы из уже принятой родительской архитектуры.
\end{nogocriterion}

\begin{protocol}
\begin{enumerate}
 \item северный и южный проекторы спинорного накрытия: \textbf{построены};
 \item экваториальный переход хопфовой линии: \textbf{$z$};
 \item коэффициентный переход: \textbf{$zq_0+1-q_0$};
 \item совпадение с $V_+$ Тома V: \textbf{точное};
 \item сопряжённая ветвь и $V_-$: \textbf{точно совпадают};
 \item намотка определителя: \textbf{$+15/-15$};
 \item граничный $K$-класс Каллиаса--Тёплица: \textbf{один и тот же};
 \item прямое унитарное равенство полных операторов: \textbf{не требуется и
 не доказано};
 \item носитель спинорного накрытия ранга два из конечного родителя:
 \textbf{не выведен};
 \item локализованные физические фермионы: \textbf{условны};
 \item следующая проверка: \textbf{родительский вывод носителя спинорного накрытия}.
\end{enumerate}
\end{protocol}

Машинный сертификат сохранён в
\path{s2t_v6_callias_toeplitz_index_comparison_gate_results.json}.